\documentclass[a4paper,10pt,landscape]{article}

\usepackage[utf8]{inputenc}
\usepackage[T1]{fontenc}
\usepackage[ngerman]{babel}
\usepackage[landscape, top=0.8cm, bottom=0.8cm, left=0.9cm, right=0.9cm]{geometry}
\usepackage{xcolor}
\usepackage{listings}
\usepackage{tcolorbox}
\usepackage{multicol}
\usepackage{fontawesome5}
\usepackage{tikz}
\usepackage{microtype}
\usepackage{lmodern}
\usepackage{parskip}

\setlength{\parskip}{2pt}
\setlength{\columnsep}{0.7cm}
\setlength{\columnseprule}{0.4pt}
\def\columnseprulecolor{\color{arduinoBlue!30}}

% ── Farben ──────────────────────────────────────────────────────────────
\definecolor{arduinoBlue}{HTML}{006D6F}
\definecolor{arduinoLightBlue}{HTML}{00A0A3}
\definecolor{goodGreen}{HTML}{2E7D32}
\definecolor{badRed}{HTML}{B71C1C}
\definecolor{codeBackground}{HTML}{F3F4F6}
\definecolor{headerBg}{HTML}{006D6F}
\definecolor{commentColor}{HTML}{5C6BC0}
\definecolor{keywordColor}{HTML}{1565C0}
\definecolor{stringColor}{HTML}{2E7D32}
\definecolor{ruleAccent}{HTML}{00ACC1}
\definecolor{ruleAccent2}{HTML}{00C174}
\definecolor{ruleAccent3}{HTML}{C1AA00}
\definecolor{ruleAccent4}{HTML}{C17700}

% ── Listings ────────────────────────────────────────────────────────────
\lstdefinestyle{cs}{
  language=C++,
  backgroundcolor=\color{codeBackground},
  basicstyle=\ttfamily\fontsize{7}{8.5}\selectfont,
  keywordstyle=\color{keywordColor}\bfseries,
  commentstyle=\color{commentColor}\itshape,
  stringstyle=\color{stringColor},
  numbers=none,
  frame=none,
  breaklines=true,
  showstringspaces=false,
  tabsize=2,
  xleftmargin=4pt,
  xrightmargin=2pt,
  morekeywords={void,int,bool,byte,float,const,String,
                pinMode,digitalWrite,digitalRead,
                analogWrite,analogRead,delay,
                Serial,begin,println,print,
                INPUT,OUTPUT,HIGH,LOW,setup,loop},
  literate={ä}{{\"a}}1 {ö}{{\"o}}1 {ü}{{\"u}}1
           {Ä}{{\"A}}1 {Ö}{{\"O}}1 {Ü}{{\"U}}1 {ß}{{\ss}}1
}
\lstset{style=cs}

% ── tcolorbox ───────────────────────────────────────────────────────────
\tcbuselibrary{breakable}

% Regelbox (ganzer Block)
\newtcolorbox{regelbox}[2]{
  colback=#1!6,
  colframe=#1,
  coltitle=white,
  fonttitle=\bfseries\fontsize{8}{9}\selectfont,
  title={#2},
  arc=3pt,
  boxrule=1pt,
  top=3pt, bottom=3pt, left=4pt, right=4pt,
  toptitle=2pt, bottomtitle=2pt
}

% Gut / Schlecht nebeneinander
\newtcolorbox{goodbox}{
  colback=goodGreen!7,
  colframe=goodGreen,
  fonttitle=\bfseries\fontsize{7}{8}\selectfont,
  title={\faCheck\enspace Richtig},
  arc=2pt, boxrule=0.8pt,
  top=2pt, bottom=2pt, left=3pt, right=3pt,
  toptitle=1pt, bottomtitle=1pt
}

\newtcolorbox{badbox}{
  colback=badRed!7,
  colframe=badRed,
  fonttitle=\bfseries\fontsize{7}{8}\selectfont,
  title={\faTimes\enspace Falsch},
  arc=2pt, boxrule=0.8pt,
  top=2pt, bottom=2pt, left=3pt, right=3pt,
  toptitle=1pt, bottomtitle=1pt
}

% ── Hilfs-Befehle ───────────────────────────────────────────────────────
\newcommand{\regelnr}[1]{%
  \begin{tikzpicture}[baseline=-0.5ex]
    \node[circle, fill=arduinoBlue, text=white,
          font=\bfseries\small, inner sep=3pt] {#1};
  \end{tikzpicture}\enspace
}

\newcommand{\sectiontitle}[2]{%
  \noindent
  {\fontsize{9}{10}\selectfont\bfseries\color{arduinoBlue}%
   \regelnr{#1}#2}\\[-2pt]
  {\color{arduinoBlue!40}\rule{\linewidth}{0.6pt}}\\[2pt]
}

% ── Header ──────────────────────────────────────────────────────────────
% 1. Logo
% 2. Überschrift
% 3. Beschreibung
% 4. Zusatztext
\newcommand{\cheatheader}[4]{%
  \noindent
  \begin{tikzpicture}
    \node[fill=headerBg, text=white, rounded corners=4pt,
          minimum width=\linewidth, minimum height=1.0cm,
          inner sep=8pt, font=\bfseries] (h)
      {\fontsize{16}{18}\selectfont {#2}\quad
       \normalfont\fontsize{10}{12}\selectfont —\quad
       {#3}};
    \node[anchor=west, xshift=8pt] at (h.west)
    		{\includegraphics[height=0.75cm]{#1}};
    \node[anchor=east, xshift=-8pt] at (h.east)
      {\color{white!80}\fontsize{8}{9}\selectfont
       {#4}};
  \end{tikzpicture}
}

\usepackage{dirtree}
\usepackage{tabularray}
\usepackage{fontawesome5}

% Dokumentenversion
\newcommand\version{V1.0}
%
% Zeilenhöhe festlegen
\newcommand\columnheight{8mm}

\definecolor{schreibflaeche}{gray}{0.95}

% ════════════════════════════════════════════════════════════════════════
\begin{document}
\pagestyle{empty}

\noindent

% Header einfügen => Definition in cheatsheet.tex
\cheatheader{../images/fablab-logo.png}{Arduino Anfängerkurs}{Spickzettel}{\today\ -\ \version}

\begin{multicols}{4}
\raggedcolumns

% Projektstruktur
\begin{regelbox}{ruleAccent}{\faFolderOpen\enspace Projektstruktur}
\dirtree{%
.1 \textbf{MeinArduinoProjekt}/.
.2 \textbf{src}/.
.3 \texttt{main.cpp}.
.2 \textbf{lib}/.
.2 \textbf{include}/.
.2 \textbf{.pio}/.
.2 \texttt{platformio.ini}.
}
\begin{tblr}{l|X}
{src} & {Programmcode} \\ \hline
{lib} & {Bibliotheken} \\ \hline
{include} & {eingebundener Code} \\ \hline
{.pio} & {PlatformIO-Dateien} \\ \hline
{platformio.ini} & {Einstellungsdatei} \\ \hline
\end{tblr}

\end{regelbox}

% Grundstruktur
\begin{regelbox}{ruleAccent2}{\faSitemap\enspace Grundstruktur des Programms}
\begin{lstlisting}
// Müssen wir bei PlatformIO immer am Anfang schreiben
#include <Arduino.h>

// Wird einmal beim Start ausgeführt
void setup() {

}

// Wird immer wieder ausgeführt
void loop() {

}
\end{lstlisting}
\end{regelbox}

% Serieller Monitor
\begin{regelbox}{ruleAccent3}{\faTv\enspace Serieller Monitor}
\begin{lstlisting}
// Seriellen Monitor starten
Serial.begin(9600);

// Ohne Zeilenumbruch ausgeben
Serial.print("Hallo");
// Mit Zeilenumbruch ausgeben
Serial.println("Welt");
\end{lstlisting}
\end{regelbox}

% Serieller Monitor
\begin{regelbox}{ruleAccent}{\faComment\enspace Kommentare}
\begin{lstlisting}
// Einfacher Kommentar

/*
Kommentar über
mehrere Zeilen
*/

delay(1000); // Kommentar

\end{lstlisting}
\end{regelbox}

% Variablen
\begin{regelbox}{ruleAccent2}{\faCode\enspace Variablen}
\begin{tblr}{l|X}
bool & Wahrheitswert \\
int & Ganzzahl \\
float & Kommazahl (16Bit) \\
double & Kommazahl (32Bit) \\
char & Zeichen \\
string & Zeichenkette \\
\end{tblr}

\begin{lstlisting}
// Variable anlegen
// Schlüsselwort Name = Wert;
int GanzeZahl = 7;

// Wert zuweisen
GanzeZahl = 5;
\end{lstlisting}

\end{regelbox}

% Bedingungen
\begin{regelbox}{ruleAccent3}{\faRandom\enspace Bedingungen}

Bedingung mit \textbf{if}. Danach in runden Klammern die Bedingung 
und in geschweiften Klammern was passieren soll.

\begin{lstlisting}
if (Bedingung) {
 // Passiert, wenn Bedingung WAHR
} else {
 // Passiert, wenn Bedingung FALSCH
}
\end{lstlisting}
\end{regelbox}

% Bedingungen
\begin{regelbox}{ruleAccent4}{\faBalanceScale\enspace Vergleichsoperatoren}
\begin{tblr}{l|X}
Operator & Funktion \\
== & ist gleich \\
!= & ist ungleich \\
< / > 	& kleiner / größer \\
<= / >= & kleiner gleich / größer gleich \\
! & nicht \\
\end{tblr}
\end{regelbox}

% Digitale Pins
\begin{regelbox}{ruleAccent}{\faBolt\enspace Digitale Pins}
\begin{lstlisting}
// Pin als Eingang festlegen
pinMode(PIN_NUMMER, INPUT);

// Pin als Eingang mit aktivem integrieren Pull-Up Widerstand
pinMode(PIN_NUMMER, INPUT_PULLUP);

// Pin als Ausgang festlegen
pinMode(PIN_NUMMER, OUTPUT);

// Zustand einlesen
digitalRead(PIN_NUMMER);

// Zustand ausgeben
digitalWrite(PIN_NUMMER, HIGH);
digitalWrite(PIN_NUMMER, LOW);
\end{lstlisting}
\end{regelbox}

% Analoge Pins
\begin{regelbox}{ruleAccent2}{\faWaveSquare\enspace Analoge Pins}

Bei analogen Pins brauchen wir kein \textbf{pinMode}.

\begin{lstlisting}
// Spannung messen
// 0V-5V => 0-1024
analogRead(PIN_NUMMER);

// Spannung ausgeben (PWM)
// wert 0-255
analogWrite(PIN_NUMMER, wert);
\end{lstlisting}

\textbf{analogWrite()} funktioniert nur bei Pins mit \textasciitilde Symbol!\\

\end{regelbox}

% Map-Funktion
\begin{regelbox}{ruleAccent2}{\faArrowRight\enspace Map Funktion}

Mit der map()-Funktion skalieren wir Werte.
z.B. wird aus 0-100 Werte von 0-255 für PWM:

\begin{lstlisting}
map(wert, vonMin, vonMax, zuMin, zuMax);

int helligkeit = map(helligkeit, 0, 100, 0, 255);
\end{lstlisting}

\end{regelbox}

% Schleifen
\begin{regelbox}{ruleAccent3}{\faSync\enspace for-Schleife}
\begin{lstlisting}
for (int i = 0; i < 10; i++) {
	// Wird 10 mal ausgeführt...
}
\end{lstlisting}
Drei Bestandteile:\\
for (Initialisierung, Bedingung, Inkrement)
\end{regelbox}

\begin{regelbox}{ruleAccent4}{\faSync\enspace while-Schleife}
\begin{lstlisting}
while (Bedingung) {
	// Wird ausgeführt solange Bedingung WAHR ist
}
\end{lstlisting}
Bedingung wie bei if/else.
\end{regelbox}

% Bibliotheken
\begin{regelbox}{ruleAccent}{\faArchive\enspace Bibliotheken}
\begin{itemize}
\item Einbinden in platformio.ini
\begin{lstlisting}
lib_deps =
	Bibliothek
	https://github.com/arduino/Bibliothek
\end{lstlisting}
\item Einbinden im Code
\begin{lstlisting}
#include "Bibliothek.h"
\end{lstlisting}
\end{itemize}
\end{regelbox}

\end{multicols}

\end{document}